\documentclass[11pt]{article}

\usepackage[T1]{fontenc}
\usepackage[utf8]{inputenc}
\usepackage{lmodern}
\usepackage[margin=1in]{geometry}
\usepackage{microtype}
\usepackage{booktabs}
\usepackage{tabularx}
\usepackage{array}
\usepackage{siunitx}
\usepackage{csquotes}
\usepackage{xurl}
\usepackage{hyperref}
\usepackage[nameinlink,noabbrev]{cleveref}
\usepackage[
  backend=biber,
  style=numeric-comp,
  sorting=none,
  maxbibnames=99,
  giveninits=true,
  doi=true,
  url=false,
  isbn=false,
  eprint=false
]{biblatex}

\addbibresource{references.bib}

\hypersetup{
  colorlinks=true,
  linkcolor=blue,
  citecolor=blue,
  urlcolor=blue,
  pdftitle={MNE-PyBlinker: a faithful Python implementation of BLINKER validated across public motor-imagery and fatigue-driving EEG datasets}
}

\DeclareFieldFormat{doi}{\mkbibacro{DOI}\addcolon\space\href{https://doi.org/#1}{#1}}
\AtEveryBibitem{
  \clearfield{issn}
  \clearfield{month}
  \clearlist{language}
}
\setlength{\emergencystretch}{3em}
\sisetup{
  group-separator = {,},
  group-minimum-digits = 4
}

\newcommand{\tool}{MNE-PyBlinker}
\newcommand{\blinker}{BLINKER}

\title{\tool: a faithful Python implementation of \blinker{} validated across public motor-imagery and fatigue-driving EEG datasets}
\author{}
\date{}

\begin{document}

\maketitle

\begin{abstract}
Automated eye-blink extraction from electroencephalography (EEG) is useful both for artifact handling and for physiological analyses, but the best known dedicated workflow, \blinker{}, remains centered on MATLAB. We report \tool, a Python implementation intended for direct use within MNE-Python pipelines, and test whether it reproduces the MATLAB reference behavior at scale rather than only approximately. We compared Python and MATLAB event tables on two public EEG resources representing distinct acquisition contexts: 74 locally staged EDF/FIF-compatible recordings from the Murat 2018 motor-imagery dataset and all 34 available EDF recordings from the SEED-VLA/VRW fatigue-driving dataset. Comparisons used harmonized blink boundaries, a \num{20}-sample tolerance window, and strict and lenient event-level metrics. Full-dataset reruns exposed two non-trivial SEED boundary failures: a truncated terminal-edge blink that survived to the final event table, and a floating-point amplitude-gate mismatch near the recording end. After correcting these cases, while rejecting a broader intermediate fix that introduced a new regression, the final Python implementation achieved exact parity with MATLAB on all 108 recordings. Across the combined benchmark, both implementations reported \num{156871} blink events, with zero Python-only and zero MATLAB-only events, and all strict and lenient macro and micro precision, recall, F1 and accuracy metrics equal to 1.0. These results establish \tool{} as a faithful, Python-native route to \blinker-style ocular-index extraction for reproducible EEG analysis.

\end{abstract}

\section{Introduction}

Spontaneous eye blinks are simultaneously a nuisance source in scalp EEG and a potentially informative physiological signal. Blink-related EEG measures have been linked to cognitive load during driving, fatigue-related monitoring, and abnormal clinical states, which means that reliable blink extraction affects both preprocessing and interpretation.\supercite{alyan2023scirep,tarafder2022,wang2018,kimura2017,mentzel2017} \blinker{} was introduced to automate the extraction of ocular indices from EEG, electrooculography and independent-component time series, making large-scale blink analysis more tractable than manual inspection.\supercite{kleifges2017}

The practical limitation is that \blinker{} remains a MATLAB workflow. That matters because contemporary open-source EEG analysis increasingly centers on Python-native tools such as MNE-Python and the wider MNE ecosystem.\supercite{gramfort2013,appelhoff2019} A Python reimplementation is only useful, however, if it preserves the reference behavior closely enough that established \blinker-based analyses can be migrated without silently changing event tables. Manual or semi-manual alternatives do not remove this problem, because annotation itself is labor-intensive and can be error-prone even when nominal agreement is high.\supercite{jansen2020}

This requirement shifts the scientific question away from detector novelty and toward implementation fidelity. Small divergences in sample indexing, terminal-edge handling, or floating-point gating can materially change blink boundaries and therefore alter downstream ocular indices. For software intended to replace a reference implementation, approximate similarity is weaker evidence than exact parity on heterogeneous public data.

Here we present \tool, a Python implementation of \blinker{} designed for use within MNE-Python workflows, and evaluate it against the MATLAB reference across two public EEG resources that cover distinct recording contexts: motor-imagery recordings from the Murat 2018 dataset\supercite{kaya2018} and fatigue-driving recordings from the public SEED-VLA and SEED-VRW datasets.\supercite{luo2024} Our aim is deliberately narrow and testable: to determine whether the Python implementation reproduces MATLAB \blinker{} outputs exactly after full-dataset reruns, and to document the edge cases that had to be resolved to reach that result.

\section{Results}

\subsection{Validation scale across two public EEG resources}

The final benchmark covered 108 recordings and \num{156871} blink events per implementation. The motor-imagery validation used 74 locally staged EDF/FIF-compatible recordings derived from the public Murat 2018 resource, whereas the fatigue-driving validation used all 34 EDF recordings available in the local SEED-VLA/VRW staging area, comprising 20 laboratory files and 14 real-world files. Event totals in \cref{tab:datasets} are counts from one implementation; Python and MATLAB totals were identical for every dataset.

\begin{table}[t]
\centering
\caption{Validation datasets and scale.}
\label{tab:datasets}
\begin{tabularx}{\linewidth}{@{}lXrr@{}}
\toprule
Dataset & Public source and local validation scope & Recordings & Blink events \\
\midrule
Murat 2018 motor-imagery EEG & Public motor-imagery dataset reported by Kaya \emph{et al.}; local validation used 74 EDF/FIF-compatible recordings processed by the fresh-comparison runner. & 74 & \num{72568} \\
SEED-VLA/VRW fatigue-driving EEG & Public laboratory and real-world fatigue-driving datasets reported by Luo \emph{et al.}; local validation used all 34 EDF recordings in the staging area (20 laboratory, 14 real-world). & 34 & \num{84303} \\
Combined & Two public EEG resources spanning motor-imagery and fatigue-driving paradigms. & 108 & \num{156871} \\
\bottomrule
\end{tabularx}
\end{table}

\subsection{Exact MATLAB--Python parity in the final retained runs}

The final retained experiment prefixes were \texttt{exp08} for Murat and \texttt{seed\_exp04} for SEED. Every recording in both reruns achieved \texttt{share\_within\_tolerance\_percent = 100.0}. At the dataset level, the Python and MATLAB pipelines produced identical blink counts, zero Python-only events, and zero MATLAB-only events. Consequently, strict and lenient macro and micro precision, recall, F1 and accuracy were all 1.0 for both datasets.

\begin{table}[t]
\centering
\caption{Final validation outcomes for the retained reruns.}
\label{tab:parity}
\footnotesize
\begin{tabular}{@{}lcccc@{}}
\toprule
Dataset & 100\% rec. & Py/MATLAB blinks & Unmatched & Strict and lenient macro/micro metrics \\
\midrule
Murat 2018 & 74/74 & \num{72568}/\num{72568} & 0/0 & All P/R/F1/Acc values = 1.0 \\
SEED-VLA/VRW & 34/34 & \num{84303}/\num{84303} & 0/0 & All P/R/F1/Acc values = 1.0 \\
Combined & 108/108 & \num{156871}/\num{156871} & 0/0 & All P/R/F1/Acc values = 1.0 \\
\bottomrule
\end{tabular}
\end{table}

\subsection{Boundary-case debugging was necessary to reach exact parity}

The perfect aggregate result in \cref{tab:parity} was not present in the first full SEED pass. Instead, complete reruns exposed two distinct boundary failures. In \texttt{real\_3}, a truncated blink extending to the last sample of the recording survived the Python pipeline and appeared as an extra terminal event. In \texttt{lab\_8}, the blink fit itself was present, but a near-end amplitude-gate decision differed by a very small floating-point margin and removed an otherwise valid event. An intermediate attempt to repair this behavior by broadening the tolerance globally restored \texttt{lab\_8} but created a new regression in \texttt{lab\_17}. The retained fix was therefore deliberately narrow: terminal-edge cleanup in the blink-geometry stage together with a boundary-only SEED rescue for the amplitude gate.

\begin{table}[t]
\centering
\caption{Boundary mismatches isolated during the full-dataset reruns.}
\label{tab:edgecases}
\small
\begin{tabularx}{\linewidth}{@{}lXX@{}}
\toprule
Recording & Failure mode & Retained resolution \\
\midrule
\texttt{real\_3} & Incomplete terminal-edge blink survived to the final event table & Explicit terminal-edge handling and cleanup \\
\texttt{lab\_8} & Near-end amplitude gate rejected a valid boundary blink & Boundary-local SEED amplitude-gate rescue \\
\texttt{lab\_17} & Regression introduced by an over-broad global tolerance change & Global tolerance change rejected \\
\bottomrule
\end{tabularx}
\end{table}

This debugging history matters because it shows that exact software migration is sensitive to edge conditions that are almost invisible in average-case summaries. The final code was kept only after targeted regression tests passed and both complete dataset reruns returned to exact parity, demonstrating that the repaired behavior fixed the boundary failures without sacrificing previously validated recordings.

\section{Discussion}

The central result of this study is implementation fidelity, not algorithmic novelty. After the final retained fixes, \tool{} reproduced MATLAB \blinker{} exactly across two public EEG resources that differ in experimental context, recording provenance and validation scale. That matters operationally: researchers can now move a mature ocular-index workflow into Python-native MNE pipelines without accepting silent changes in blink-event tables as the cost of migration.\supercite{gramfort2013,appelhoff2019}

The work also shows why a faithful port must be validated at the level of full-dataset reruns rather than spot checks. The remaining mismatches were not broad failures of the detector, but narrow differences in terminal-edge geometry and boundary-local floating-point behavior. Those are precisely the classes of software discrepancies most likely to survive superficial regression testing while still altering downstream analyses.

\paragraph{Self-critique and limitations.}
The strongest claim supported here is exact parity with the MATLAB reference, not superiority over alternative blink detectors and not independent confirmation that every blink boundary is physiologically optimal. First, the reference standard was MATLAB \blinker{} itself rather than a newly curated human-annotation gold standard, so any systematic bias present in the legacy implementation is inherited by this validation. Second, the evaluation spans only two public resources and the Murat benchmark used locally staged EDF/FIF-compatible derivatives rather than every original raw distribution format. Third, the final SEED result required a boundary-local numerical rescue in the validation harness after a floating-point mismatch was isolated near the recording end; this should be read as careful software replication, not as evidence that one globally tuned parameter set is universally sufficient. Fourth, we did not benchmark runtime, memory use, or the downstream stability of higher-order ocular indices beyond event boundaries.

Those limitations define the next steps. Independent human-annotation benchmarking, broader validation across additional device families and patient cohorts, direct assessment of downstream ocular-index reproducibility, and public archival release of a versioned software snapshot would all strengthen the package further. Even with those caveats, the present result is useful: a previously MATLAB-bound workflow can now be executed in Python with auditable, recording-level evidence that the migrated implementation preserves the legacy event tables exactly on the tested public datasets.

\section{Methods}

\subsection{Study design}

The study was designed as a software-fidelity benchmark. MATLAB \blinker{} outputs were treated as the reference implementation, and \tool{} outputs were compared recording by recording against those reference event tables. A code change was retained only if it resolved the targeted mismatch and preserved perfect results in complete reruns of both the Murat and SEED validation suites.

\subsection{Datasets}

The first benchmark used the public Murat 2018 motor-imagery EEG resource reported by Kaya \emph{et al.}\supercite{kaya2018} The source publication describes 75 long recording sessions. The local validation harness processed the 74 recordings for which EDF or FIF-compatible inputs were available to the fresh-comparison runner. The second benchmark used the public SEED-VLA and SEED-VRW fatigue-driving EEG datasets described by Luo \emph{et al.}\supercite{luo2024} The local staging area contained 34 EDF recordings, split into 20 laboratory recordings and 14 real-world recordings. No new human participant data were collected for this work.

\subsection{Python and MATLAB pipelines}

Python detections were generated through \tool{} using MNE-Python I/O.\supercite{gramfort2013} Raw EDF or FIF recordings were loaded with MNE, filtered from 1 to \SI{20}{Hz}, and processed at the native sampling rate after rounding the sampling frequency to the nearest integer for resampling. The Murat benchmark used the explicit legacy-default blink-parameter dictionary preserved in the repository: \texttt{std\_threshold = 1.5}, \texttt{min\_event\_len = 0.05}, \texttt{min\_event\_sep = 0.05}, \texttt{base\_fraction = 0.1}, correlation thresholds of \texttt{0.98/0.90/0.95} for top, bottom and middle segments, \texttt{shut\_amp\_fraction = 0.9}, blink amplitude range \texttt{[3, 50]}, \texttt{good\_ratio\_threshold = 0.7}, \texttt{min\_good\_blinks = 10}, \texttt{correlation\_threshold = 0.98}, \texttt{p\_avr\_threshold = 3}, and \texttt{z\_thresholds = [[0.9, 0.98], [2.0, 5.0]]}. The final SEED rerun used the same legacy defaults together with a boundary-only numerical safeguard, \texttt{amplitude\_gate\_tolerance = 5e-8} and \texttt{amplitude\_gate\_end\_window\_seconds = 30.0}, introduced after the near-end floating-point mismatch had been isolated.

MATLAB reference outputs were produced with the EEGLAB \blinker{} plugin (version 1.2.0) through the repository's validation harness. For the SEED rerun this produced MATLAB pickles adjacent to each EDF file; for Murat the reference pickles were the existing per-recording \texttt{blinker\_results.pkl} artifacts used by the fresh-comparison runner.

\subsection{Event-table harmonization and comparison metrics}

The Python implementation stores blink boundaries as zero-based \texttt{left\_zero} and \texttt{right\_zero} indices. Before comparison, these columns were renamed to \texttt{start\_blink} and \texttt{end\_blink} and shifted by +1 sample so that the Python boundaries matched the one-based MATLAB convention. Comparisons then used a \num{20}-sample tolerance window.

Per-recording summaries were built from the pairwise difference table produced by the evaluation utilities. Events classified as \texttt{share\_within\_tolerance} counted as strict true positives. Events classified as \texttt{matches\_within\_tolerance} satisfied the boundary window but failed at least one overlap or amplitude check; these were excluded from the strict count and included in the lenient count. False positives and false negatives were \texttt{detected\_only} and \texttt{ground\_truth\_only}, respectively. Macro metrics were obtained by averaging per-recording values, whereas micro metrics were computed from pooled totals across recordings.

\subsection{Rerun discipline and artifact tracking}

Long reruns were executed under explicit experiment prefixes so that each logic revision produced a distinct artifact set rather than overwriting previous results. The final retained prefixes were \texttt{exp08} for Murat and \texttt{seed\_exp04} for SEED. The validation harness wrote per-recording pickle outputs, rolling logs, live-status files for long jobs, and final summary CSV and overall JSON files. This file-based audit trail was treated as part of the methods because it allowed every retained change to be checked against fixed artifacts after the run completed.

\section*{Data availability}

No new datasets were generated in this study. The Murat motor-imagery EEG data analysed here are publicly distributed through Figshare and described by Kaya \emph{et al.}\supercite{kaya2018} The fatigue-driving data are the public SEED-VLA and SEED-VRW datasets described by Luo \emph{et al.}\supercite{luo2024} and distributed through the SEED portal at Shanghai Jiao Tong University. The present study used local EDF/FIF-compatible staged copies of those public resources for direct MATLAB--Python comparison.

\section*{Code availability}

The source tree used for the study is included with the submission as the repository \path{blinker_pyblinker_validation}. The exact retained experiment artifacts are \path{reports/validation/exp08_top74_summary.csv}, \path{reports/validation/exp08_top74_overall.json}, \path{reports/validation/seed_exp04_seed_summary.csv}, and \path{reports/validation/seed_exp04_seed_overall.json}. Investigation notes documenting the retained fix are archived in \path{issue.md} and \path{docs/findings/034_seed_murat_edge_fix_rerun.md}. The repository also preserves the validation scripts that generated these artifacts, including the fresh Murat sweep runner and the full SEED pipeline.

\section*{Ethics statement}

This study re-analysed publicly released, de-identified EEG datasets. No new participant recruitment, intervention or contact occurred as part of this work. Ethics approval and informed-consent procedures for the original data collections are described in the source publications.\supercite{kaya2018,luo2024}


\printbibliography[title={References}]

\end{document}
